\section{序論}\label{sec:intro}

\subsection{背景}
U$(N)$ Chern--Simons 理論の彩色 HOMFLY-PT 不変量は、resolved conifold 上の開弦の位相的弦と large $N$ で双対であり、Ooguri--Vafa の再定式化によって整数の BPS 不変量（LMOV 不変量）で書ける。
SO$(N)$・Sp$(N)$ の場合は、Calabi--Yau の orientifold の位相的弦が対応する。Mariño \cite{Marino09} は、SO$(N)$ の彩色 Kauffman 不変量から、U$(N)$ の composite 表現の寄与の半分を引いた残りが、クロスキャップ $c = 1, 2$ の整数 BPS 不変量で書けることを予想し、多くの例で確かめた（Bouchard--Florea--Mariño \cite{BFM04,BFM05}、Nawata--Ramadevi--Zodinmawia \cite{NRZ13} なども参照）。

refined の側では、Aganagic--Shakirov \cite{AgSh11} が refined Chern--Simons 理論を定式化した。著者と Nawata \cite{KN17} は、U$(N)$ の refined Chern--Simons 不変量が resolved conifold 上の refined BPS 不変量を与え、トーラス結び目で正の整数になることを確かめた（refined large $N$ 双対）。
Aganagic--Schaeffer \cite{AgSc12} は、orientifold の refined 位相的弦を M 理論の指数として定義した。そのうえで SO$(2N)$ の refined Chern--Simons 不変量を計算し、大きな $N$ の閉弦の自由エネルギーが第二量子化の整数性を満たすことを示した。
ただし開弦（結び目）については、refined の BPS 不変量を取り出す式も、その整数性・正値性の検証も無かった。
著者は以前（2018 年）に、Mariño の構造を素朴に refined にする計算を行ったが、整数性が得られなかった。

\subsection{本論文の問い}
本論文では次の問いを扱う。
\begin{enumerate}
\item SO$(2N)$ の refined Chern--Simons 不変量から、refined のクロスキャップの BPS 不変量を取り出す式はどう書けるか。
\item その整数性はどの範囲で成り立つか。
\item refined で正値性は見えるか。見えないなら、何が障害か。
\end{enumerate}

\subsection{結果の概要}
\begin{itemize}
\item \textbf{定式化（§\ref{sec:prelim}）}：Mariño 型の bosonic な核を素朴に refined にすると、unknot の 2 箱で既に $U = t^N$ の極が打ち消さない。代わりに双対 Cauchy 核（t-brane）を使うと、unknot の積表示が自動的に成り立つ。
規格化は、U$(n)$ の不変量が行列式と荷電共役で不変になる「一様な規格化」をとる。reformulation には Mariño の $M(q^{-1})$ を使う。
\item \textbf{構造（§\ref{sec:index}）}：orientifold の M 理論の指数から、全 $d$ の構造 $g = \tfrac12 E/(1+t) + C$ を導く（命題 \ref{prop:alld}）。$E$ と $C$ の分離は指数からは決まらない（命題 \ref{prop:EC}）。
この構造は、Witten \cite{Witten11} の Chern--Simons 分配関数の trace 表示の SO$(2N)$ 版として読める（命題 \ref{prop:trace}）。境界ゼロモードへの $\tilde I$ の作用から、$E$ の境界因子が $z_q^{n-1}(z_q^+)^n$ になることを導き、Hopf 絡み目・結び目・framed unknot で確かめた（命題 \ref{prop:boundary}）。
閉弦では、SO$(2N)$ の分配関数の積表示が $N$ の 1 次式を除いて Aganagic--Schaeffer の式と一致することを確かめ、Avetisyan--Mkrtchyan \cite{AM22} の「Klein 瓶」の寄与がこの分解の並べ替えであることを示す（命題 \ref{prop:AM}）。
\item \textbf{unknot と Hopf 絡み目（§\ref{sec:hopf}）}：整数性が成り立つ。Kähler 変数は $\lambda \to \tlam = \lambda(q/t)^{1/2}$ とずれ、$E_{22}, E_{24}$ は一意に決まる。
\item \textbf{トーラス結び目の [1]（§\ref{sec:torus}）}：$T(2,2m+1)$（$m \le 5$）で、$c = 1$ 部分の $G$ 展開と $c = 2$ 部分の閉じた形が得られ、符号がそろう判定式が成り立つ。
\item \textbf{一般の $N$ の多箱（§\ref{sec:generalN}）}：DAHA-Jones 多項式を一般の $t$ で計算し、refined SO$(2N)$ の trefoil の 2 箱・3 箱と $T(2,5)$ の 2 箱を一般の $N$ で求めた。
クロスキャップ部分は整数で、$(1+t)$ の極は $p_2$ 方向にだけ現れる。refined だけの部分は普遍因子 $(\bt^2 - 1)(\ba^2 - \bq^2)$ で割り切れる。
\item \textbf{refined Sp$(2N)$（§\ref{sec:sp}）}：C$_N$ 型の DAHA を実装し、Avetisyan--Mkrtchyan の refined Sp$(2N)$ の trefoil・$T(2,5)$ の [1] と trefoil の 2 箱を一般の $N$ で求めた。[1] では $g^{\mathrm{Sp}}(\ba) = -g^{\SO}(-\ba)$、2 箱では $W^{\mathrm{Sp}}_R(\ba) = (-1)^{|R|} W^{\SO}_R(-\ba)$ で、クロスキャップ部分の差は閉じた形 $\pm\psi^2[\delta]$ に尽きる。
\item \textbf{正値性（§\ref{sec:positivity}）}：多箱では符号がそろわない。障害は、凍結した状態の多重項の基底と、$E$・$C$ の分離にある。$E$ と $C$ の分離は指数からは決まらず、composite の DAHA の値からも状態の数は得られなかった。framed unknot では $E$ が framing によらず一意に決まる。多箱の $c = 1$ の $\sigma$ 非対称性は、$c$ の偶奇の規則（hatted では $\beta \equiv |R| - 1$）と、refined に固有の $\sigma$ 奇部分に分かれる。
\end{itemize}

\begin{table}[t]
\centering
\small
\caption{本論文の検証の範囲。「整数」は整数係数の Laurent 多項式（$|R| = 2$ は $2g$ が整数で $g_{(2)} + g_{(1,1)}$ が整数）。「普遍因子」は §\ref{sec:universal} の割り切れ。}\label{tab:summary}
\begin{tabular}{p{3.6cm}p{2.6cm}p{2.4cm}p{5.4cm}}
\toprule
対象 & 整数 & 普遍因子 & 正値性（符号のそろい方） \\
\midrule
unknot、Hopf（$a, b \le 4$） & 成立 & 成立（$\hat C_{R,S}$ の 10 成分） & $\hat g_{R,[1]}$ はそろう。$\hat C_{R,S}$ の refined だけの部分は 8/10 \\
$T(2,2m+1)$ の [1]（$m \le 5$） & 成立 & $c = 2$ 部分は $m \le 5$、全体は $m \le 2$ & 判定式が成立（商の符号は $m \le 2$ で確認） \\
$T(2,2m+1)$ の固定 $N$（$N = 2, 3$、$|R| \le 3$、$m \le 3$） & 全方向の hatted で成立 & — & 固定 $N$ では判定できない \\
trefoil の 2 箱（大きな $N$） & 成立 & 成立 & 商のブロックの 5/10 \\
$T(2,5)$ の 2 箱（大きな $N$） & 成立 & 成立 & 6/18 \\
trefoil の 3 箱（大きな $N$） & 成立（$(1+t)$ の極なし、hatted も整数） & 成立 & 2/15, 3/14, 2/13 \\
refined Sp$(2N)$：$T(2,3)$, $T(2,5)$ の [1]、trefoil の 2 箱 & 成立（SO と同じ構造） & — & SO との関係 $g^{\mathrm{Sp}}(\ba) = -g^{\SO}(-\ba)$（[1]）、$W_R$ の $(-1)^{|R|}$（2 箱） \\
\bottomrule
\end{tabular}
\end{table}

\subsection{主張の分類}
本論文の主張は、解析的に導いたもの（定理・命題・補題）、計算機で確かめたもの（計算結果。範囲を明記する）、予想、仮説に分けて書く。
文献から直接取れなかった規則は「仮説」と書き、どの検証で裏付けたかを添える。
計算のコードとデータ、自動テストの対応は付録 \ref{app:code} に置く。
